% Poisson MLE, full worked solution (four steps)
% — source recovered from wk02 deck notes, 2026-07-22
\documentclass[border=3bp]{standalone}

% decent font to pair with Source Sans 3
\usepackage{tgpagella}  

% paragraph formatting
\setlength{\parindent}{0pt}
\setlength{\parskip}{0.75em}

% color definitions
\usepackage[usenames,dvipsnames]{xcolor}
\definecolor{red}{HTML}{e41a1c}
\definecolor{blue}{HTML}{377eb8}
\definecolor{green}{HTML}{4daf4a}

% packages
\usepackage{amssymb}
\usepackage{amsmath}
\usepackage{mathrsfs}
\usepackage[utf8]{inputenc}


\begin{document}

\begin{minipage}{4.5in}
\textbf{Likelihood.}
\[
L(\lambda)=\prod_{i=1}^{N}\frac{e^{-\lambda}\lambda^{y_i}}{y_i!}
= e^{-N\lambda}\,\lambda^{\sum_{i=1}^{N}y_i}\,\prod_{i=1}^{N}\frac{1}{y_i!}.
\]

\begin{enumerate}\setlength{\itemsep}{0.25em}

\item \textbf{Take logs.} (Use $\log(\prod a_i)=\sum\log a_i$, $\log(ab)=\log a+\log b$, and $\log(x^c)=c\log x$.)
\begingroup\setlength{\jot}{1.5pt}
\[
\begin{aligned}
\log L(\lambda)
&= \sum_{i=1}^{N}\left[-\lambda + y_i\log\lambda - \log(y_i!)\right] \\
&= -N\lambda + \Big(\sum_{i=1}^{N}y_i\Big)\log\lambda - \sum_{i=1}^{N}\log(y_i!).
\end{aligned}
\]
\endgroup

\item \textbf{Differentiate $\log L(\lambda)$ with respect to $\lambda$.}
\begingroup\setlength{\jot}{1.5pt}
\[
\begin{aligned}
\frac{d\,\log L(\lambda)}{d\lambda}
&= \frac{d}{d\lambda}\!\left[-N\lambda\right]
  + \frac{d}{d\lambda}\!\left[\Big(\sum_{i=1}^{N}y_i\Big)\log\lambda\right]
  - \frac{d}{d\lambda}\!\left[\sum_{i=1}^{N}\log(y_i!)\right] \\
&= -N \;+\; \Big(\sum_{i=1}^{N}y_i\Big)\frac{1}{\lambda} \;+\; 0 \\
&= -\,N \;+\; \frac{\sum_{i=1}^{N}y_i}{\lambda}.
\end{aligned}
\]
\endgroup

\item \textbf{Set the derivative to zero at $\lambda=\hat{\lambda}$ and solve.}
\[
\left.\frac{d\,\log L(\lambda)}{d\lambda}\right|_{\lambda=\hat{\lambda}}
= -\,N + \frac{\sum_{i=1}^{N}y_i}{\hat{\lambda}} = 0.
\]
\begingroup\setlength{\jot}{1.5pt}
\[
\begin{aligned}
\frac{\sum_{i=1}^{N}y_i}{\hat{\lambda}} &= N
\;\Longleftrightarrow\;
\sum_{i=1}^{N}y_i = N\hat{\lambda}
\;\Longrightarrow\;
\hat{\lambda} = \frac{1}{N}\sum_{i=1}^{N}y_i
= \overline{y}.
\end{aligned}
\]
\endgroup

\item \textbf{Second-order condition (concavity).} Need $\displaystyle \frac{d^{2}\!\log L(\lambda)}{d\lambda^{2}}<0$.
\[
\frac{d^{2}\!\log L(\lambda)}{d\lambda^{2}}
= -\,\frac{\sum_{i=1}^{N}y_i}{\lambda^{2}}
\;<\;0 \quad (\lambda>0\ \text{and}\ \textstyle\sum_{i=1}^{N}y_i>0).
\]
Thus $\log L(\lambda)$ is strictly concave and the critical point $\hat{\lambda}=\overline{y}$ is the unique global maximizer.
\end{enumerate}
\end{minipage}

\end{document}
